% !TeX root = ../ResponsetoReviewers.tex
\documentclass[../ResponsetoReviewers.tex]{subfiles}
\begin{document}

\begin{reviewersection}

\begin{reviewercomment}
    The authors developed a self-consistent field theory for an AB diblock copolymer with a sliding A-homopolymer on the A-block.
    This architecture is very interesting because it is changeable or even responsive to local environment.
    As this manuscript was previously submitted to this journal and was reviewed by me, the authors have made significant revisions and have addressed most of my comments.
    I just have some minor comments for them.
\end{reviewercomment}


\begin{qa}
    This kind of architecture with a sliding block is responsive to the local self-assembly environment, and therefore it should be more adaptable to nonuniform domain thickness or interface curvature.
    Is it possible to take advantage of the changeable architecture to stabilize novel ordered phases?
\end{qa}
\begin{reply}
    Response to be added.
\end{reply}


\begin{qa}
    In the strong segregation limit, the sliding block tends to move to the end of the A block, becoming close to a linear AB deblock copolymer.
    Accordingly, the length ratio of the sliding block, $\phi_{\mathrm a}$, should have little effect on the phase boundaries.
    When I look at the phase diagrams in Figure 3, there exist some noticeable differences between the phase boundaries of this special copolymer and AB diblock copolymer.
    Do the differences taper off as increasing $\chi N$?
\end{qa}
\begin{reply}
    Response to be added.
\end{reply}

\end{reviewersection}

\end{document}
